\documentclass[10pt]{extarticle}
\usepackage[a4paper,margin=1.5cm]{geometry}
\usepackage[T1]{fontenc}
\usepackage{amsmath,amssymb}
\usepackage{graphicx}
\usepackage[export]{adjustbox}
\usepackage{multicol}
\usepackage{enumitem}
\usepackage{booktabs,array,tabularx}
\setlength{\columnsep}{0.7cm}
\setlength{\parindent}{0pt}
\setlength{\parskip}{3pt}
\pagestyle{empty}
\makeatletter
% figure registry (figdefs.tex is generated by the extract step)
\newcommand{\qfigdef}[4]{\@namedef{qf@file@#1}{#2}\@namedef{qf@w@#1}{#3}\@namedef{qf@wide@#1}{#4}}
\newcommand{\qf@check}[1]{\@ifundefined{qf@file@#1}{\PackageError{score2}{Unknown figure id #1}{}}{}}
\newcommand{\qf@img}[4][]{\includegraphics[width=#2\dimexpr\@nameuse{qf@w@#4}pt\relax,max width=#3,max height=0.45\textheight,#1]{\@nameuse{qf@file@#4}}}
% Full-width questions: finalize wraps every question that contains a wide figure in \qfullbegin...\qfullend,
% so the WHOLE question (stem, figures, options) is set across the page in reading order, then the two columns
% resume. The question's head - its number, stem and first figure (\qfullsplit marks the end of the head) -
% is kept on one page, so a number is never left at a page foot with its figure overleaf; the rest may break
% across pages normally, so no big gaps are left. Without a \qfullsplit, the whole question is kept together
% if it is under 45% of a page. Never write these yourself.
\newif\ifqf@full
\newif\ifqf@boxed
\newsavebox{\qf@box}
\newcommand{\qf@open}{\global\qf@fulltrue\par\global\qf@boxedtrue
  \setbox\qf@box\vbox\bgroup\hsize\textwidth\linewidth\textwidth}
\newcommand{\qf@emit}[1]{\par\egroup\global\qf@boxedfalse
  \ifdim\dimexpr\ht\qf@box+\dp\qf@box\relax<#1\textheight\noindent\usebox{\qf@box}\par\else\unvbox\qf@box\fi}
\newcommand{\qfullsplit}{\ifqf@boxed\qf@emit{0.85}\fi}
\newcommand{\qf@close}{\ifqf@boxed\qf@emit{0.45}\fi\global\qf@fullfalse}
\newcommand{\qfullbegin}{\end{multicols}\qf@open}
\newcommand{\qfullend}{\qf@close\begin{multicols}{2}}
\newcommand{\qfullstart}{\qf@open}
\newcommand{\qfulldone}{\qf@close}
% \qfig[scale]{id}: block figure, centred, directly where it is written (size set by the tool).
% Fallback if a wide figure is not inside a full-width question: break out of the columns in place - no floats.
\newcommand{\qfig}[2][1]{\qf@check{#2}%
  \ifqf@full\qf@blk{#1}{#2}\else
  \ifnum\@nameuse{qf@wide@#2}=1 \expandafter\@firstoftwo\else\expandafter\@secondoftwo\fi
  {\end{multicols}\begin{center}\qf@img{#1}{\textwidth}{#2}\end{center}\begin{multicols}{2}}%
  {\qf@blk{#1}{#2}}\fi}
\newcommand{\qf@blk}[2]{\par\smallskip{\centering\qf@img{#1}{\linewidth}{#2}\par}\smallskip}
% \qfigi[scale]{id}: inline figure (inside option lists / tables / minipages); never breaks columns.
\newcommand{\qfigi}[2][1]{\qf@check{#2}\qf@img[valign=c]{#1}{\linewidth}{#2}}
\makeatother
% structure macros - lint reads these to check completeness and figure placement
\newcommand{\papersec}[1]{\par\bigskip{\centering\Large\bfseries PAPER-#1\par}\medskip}
\newcommand{\partsec}[2]{\par\medskip\noindent{\large\bfseries PART-#1: #2}\par\smallskip}
\newcommand{\sect}[1]{\par\smallskip\noindent\textbf{#1}\par}
\newcommand{\qn}[1]{\par\smallskip\noindent\textbf{#1.}~}
\newcommand{\paraq}[2]{\par\smallskip\noindent\textbf{Paragraph for Questions #1 and #2}\par}
\newenvironment{qoptions}{\begin{itemize}[leftmargin=1.8em,itemsep=1pt,topsep=2pt,label={}]}{\end{itemize}}
\newcommand{\opt}[2]{\item[(#1)] #2}
% figure registry (generated by score2_tools.py finalize)
\qfigdef{p1_phy_s0_q01_1}{SRGBT - Test 12 (10 Oct 2021)-img/p1_phy_s0_q01_1.png}{46.7}{0}
\qfigdef{p1_phy_s0_q05_2}{SRGBT - Test 12 (10 Oct 2021)-img/p1_phy_s0_q05_2.png}{122.5}{0}
\qfigdef{p1_phy_s0_q05_1}{SRGBT - Test 12 (10 Oct 2021)-img/p1_phy_s0_q05_1.png}{118.3}{0}
\qfigdef{p1_phy_s0_q06_1}{SRGBT - Test 12 (10 Oct 2021)-img/p1_phy_s0_q06_1.png}{87.6}{0}
\qfigdef{p1_phy_s0_q07_1}{SRGBT - Test 12 (10 Oct 2021)-img/p1_phy_s0_q07_1.png}{97.9}{0}
\qfigdef{p1_phy_s0_q09_1}{SRGBT - Test 12 (10 Oct 2021)-img/p1_phy_s0_q09_1.png}{58.0}{0}
\qfigdef{p1_phy_s0_q10_1}{SRGBT - Test 12 (10 Oct 2021)-img/p1_phy_s0_q10_1.png}{133.5}{0}
\qfigdef{p1_phy_s0_q16_1}{SRGBT - Test 12 (10 Oct 2021)-img/p1_phy_s0_q16_1.png}{125.1}{0}
\qfigdef{p1_chem_s0_q24_1}{SRGBT - Test 12 (10 Oct 2021)-img/p1_chem_s0_q24_1.png}{56.9}{0}
\qfigdef{p1_chem_s0_q25_1}{SRGBT - Test 12 (10 Oct 2021)-img/p1_chem_s0_q25_1.png}{228.5}{0}
\qfigdef{p1_chem_s0_q32_1}{SRGBT - Test 12 (10 Oct 2021)-img/p1_chem_s0_q32_1.png}{170.7}{0}
\raggedcolumns
\makeatletter\@beginparpenalty=10000\makeatother
\newcolumntype{L}[1]{>{\raggedright\arraybackslash}p{#1}}
\begin{document}
\begin{center}{\Large\bfseries SRGBT --- Test 12 (10 Oct 2021)}\end{center}
\begin{multicols}{2}
\raggedcolumns
\papersec{1}
\partsec{1}{PHYSICS}
\qn{1} The system shown in the adjacent figure is in equilibrium. The mass of the container with liquid is M, density of liquid in the container is $\rho$ and volume of the block is V. If the container is now displaced downwards through a distance x$_0$ and released such that the block remains well inside the liquid, then during subsequent motion\par
\qfig{p1_phy_s0_q01_1}
\begin{qoptions}
\opt{A}{time period of SHM of the container will be $2\pi\sqrt{\dfrac{\mathrm M}{\mathrm k}}$}
\opt{B}{time period of SHM of the container will be $2\pi\sqrt{\dfrac{(\mathrm M+\rho\mathrm v)}{\mathrm k}}$}
\opt{C}{amplitude of SHM of the container will be greater than x$_0$}
\opt{D}{container will not perform SHM}
\end{qoptions}
\qn{2} A particle of mass m is performing simple harmonic motion along line AB with amplitude 2a with centre of oscillation as O. At t = 0 particle is at point C (OC = a) and is moving towards B with velocity v = a$\sqrt3$ m/s. The equation of motion can be given by\par
\qfig{p1_phy_s0_q05_2}
\begin{qoptions}
\opt{A}{x = 2a (sin t + $\sqrt3$ cos t)}
\opt{B}{x = 2a ( $\sqrt3$ sin t + cos t)}
\opt{C}{x = a (sin t + $\sqrt3$ cos t)}
\opt{D}{x = a ( $\sqrt3$ sin t + cos t)}
\end{qoptions}
\qn{3} 200 students are talking in an examination in a room, and the sounds of pens scratching on paper, sighs, groans, and muttered imprecations has created a more or less continuous sound level of this noise of 60dB. Assuming each student contributes equally to this noise and nothing else changes or adds to it, what will the sound level in the room be when only 50 students are left ? ($\log_{10}2=0.3$)
\begin{qoptions}
\opt{A}{50 dB}
\opt{B}{15 dB}
\opt{C}{66 dB}
\opt{D}{54 dB}
\end{qoptions}
\qn{4} A sonometer wire vibrates with a frequency f. It is replaced by another wire of three times the diameter if tension and other parameters remain the same, the frequency of vibration of the wire will be-
\begin{qoptions}
\opt{A}{$\dfrac{\mathrm f}9$}
\opt{B}{$\dfrac{\mathrm f}3$}
\opt{C}{3f}
\opt{D}{9f}
\end{qoptions}
\qn{5} Here snap shot of a progressive wave at $\mathrm t=\dfrac{\mathrm T}6$ is given with time period = T. Then the equation of the wave will be if wave is going in positive x-direction and if wave is going in negative x-direction respectively. $\left(\text{Here, T}=\dfrac{2\pi}{\Omega}\right)$ (symbol have their usual meaning)\par
\qfig{p1_phy_s0_q05_1}
\begin{qoptions}
\opt{A}{$y_1=\mathrm A\sin\left(\mathrm{kx}-\omega\mathrm t+\dfrac\pi2\right)$, $y_2=\mathrm A\sin\left(\mathrm{kx}+\omega\mathrm t-\dfrac\pi6\right)$}
\opt{B}{$y_1=\mathrm A\cos(\mathrm{kx}-\omega\mathrm t)$, $y_2=\mathrm A\sin(\mathrm{kx}+\omega\mathrm t)$}
\opt{C}{$y_1=\mathrm A\sin(\mathrm{kx}-\omega\mathrm t)$, $y_2=\mathrm A\sin\left(\mathrm{kx}+\omega\mathrm t-\dfrac\pi6\right)$}
\opt{D}{$y_1=\mathrm A\sin\left(\mathrm{kx}-\omega\mathrm t+\dfrac\pi2\right)$, $y_2=\mathrm A\sin\left(\mathrm{kx}+\omega\mathrm t+\dfrac\pi6\right)$}
\end{qoptions}
\qn{6} You have a microwave generator that can produce microwaves at any frequency between 1GHz and 10 GHz. The microwave radiation enters a 10-cm-long cylinder with reflective end caps, as shown in the Figure. What frequency of the microwave generator, in GHz, will produce the lowest-order standing wave with an antinode (maximum) in the center of the cavity? Note that with reflectors at both ends, the electromagnetic standing wave acts just like the standing wave on a string that is tied at both ends. Take velocity of waves as 3 $\times$ 10$^8$ m/s.\par
\qfig{p1_phy_s0_q06_1}
\begin{qoptions}
\opt{A}{4.5}
\opt{B}{3.0}
\opt{C}{1.5}
\opt{D}{7.5}
\end{qoptions}
\qn{7} A particle is suspended by two ideal strings as shown in the figure. Now mass m is given a small displacement perpendicular to the plane of triangle formed. Choose the correct statement(s).\par
\qfig{p1_phy_s0_q07_1}
\begin{qoptions}
\opt{A}{The period of oscillation of the system is $2\pi\sqrt{\dfrac{3\sqrt3\,\ell}{\mathrm g}}$.}
\opt{B}{The period of oscillation of the system is $2\pi\sqrt{\dfrac{3\ell}{\mathrm g}}$.}
\opt{C}{The period of oscillation of the system is independent of M.}
\opt{D}{If the distance between the suspension points was kept constant and the lengths of the strings were quadrupled then the period of the system will double.}
\end{qoptions}
\qn{8} Two particles A and B are performing SHM along x and y-axis respectively with equal amplitude and frequency of 2 cm and 1 Hz respectively. Equilibrium positions of the particles A and B are at the co-ordinates (3, 0) and (0, 4) respectively. At t = 0, B is at its equilibrium position and moving towards the origin, while A is nearest to the origin and moving away from the origin. Choose the correct option(s) :
\begin{qoptions}
\opt{A}{The distance between A and B is 3 cm at t = t$_0$ sec. It will again be 3 cm at t = (t$_0$ + 1) sec.}
\opt{B}{The distance between A and B can never be 5 cm.}
\opt{C}{The maximum magnitude of force acting on each particle during the motion is same.}
\opt{D}{Kinetic energy of the system (A and B) is conserved}
\end{qoptions}
\qn{9} A torsional pendulum consists of two identical rods each of mass M \& length $\ell$ joined perpendicular to each other in horizontal plane and is hanged through the wire having torsional constant C. The system performs simple harmonic motion $\theta=\mathrm A\sin(\omega\mathrm t+\phi)$ such that initially angular displacement of the system is equal to half of amplitude and its angular velocity is $\omega_0$ and is decreasing, then :-\par
\qfig{p1_phy_s0_q09_1}
\begin{qoptions}
\opt{A}{Time period of pendulum is $\mathrm T=2\pi\sqrt{\dfrac{\mathrm M\ell^2}{6\mathrm C}}$}
\opt{B}{Amplitude of oscillation is $\dfrac{\omega_0\ell}3\sqrt{\dfrac{2\mathrm M}{\mathrm C}}$}
\opt{C}{Initial phase of SHM is $\dfrac\pi6$ or $\dfrac{7\pi}6$}
\opt{D}{Energy stored in wire at t = 0 is $\dfrac{\mathrm M\ell^2\omega_0^2}{36}$}
\end{qoptions}
\qn{10} Two identical straight wires PQ and RS each of mass m and length $\ell$ can move smoothly on a fixed rectangular frame. Two thin films of a liquid of surface tension T are formed between each wire and the frame. The two wires are connected by a massless spring of stiffness k and intially in a natural length position and released then choose the correct option(s).\par
\qfig{p1_phy_s0_q10_1}
\begin{qoptions}
\opt{A}{Maximum elongation of spring, $\Delta x_\mathrm m=\dfrac{2\mathrm T\ell}{\mathrm k}$}
\opt{B}{Maximum elongation of spring, $\Delta x_\mathrm m=\dfrac{4\mathrm T\ell}{\mathrm k}$}
\opt{C}{Each wire executes SHM with the time period, $\mathrm T_0=2\pi\sqrt{\dfrac{\mathrm m}{2\mathrm k}}$}
\opt{D}{Each wire executes SHM with the time period, $\mathrm T_0=2\pi\sqrt{\dfrac{2\mathrm m}{\mathrm k}}$}
\end{qoptions}
\qn{11} The first overtone of an open organ pipe (of length L$_0$) beats with the first overtone of a closed organ Pipe (of length L$_\mathrm c$) with a beat frequency of 10 Hz. The fundamental frequency of the closed pipe is 110 Hz. If the speed of sound is 330 m/s. Then
\begin{qoptions}
\opt{A}{L$_\mathrm c$ = 0.75 m}
\opt{B}{L$_\mathrm o=\left(\dfrac{33}{34}\right)$m or $\left(\dfrac{33}{35}\right)$m}
\opt{C}{L$_\mathrm o=\left(\dfrac{33}{32}\right)$m or 1 m}
\opt{D}{L$_\mathrm o=\left(\dfrac{33}{32}\right)$m or $\left(\dfrac{33}{34}\right)$m}
\end{qoptions}
\qn{12} A vibrator makes 150 cm of a string to vibrate in 6 loops in the longitudinal arrangement when it is stretched by 15 g. Then The entire length of the string is then weighed and is found to weigh 500 milligram.
\begin{qoptions}
\opt{A}{frequency of the vibrator is 84 Hz}
\opt{B}{frequency of the vibrator is 42.4 Hz.}
\opt{C}{distance between two nodes is 25 cm}
\opt{D}{distance between two nodes is 33 cm}
\end{qoptions}
\qn{13} The wave function for the wave pulse is $y(x,\,t)=\dfrac{0.1\mathrm C^3}{\mathrm C^2+(x-\mathrm{vt})^2}$, with C = 4 cm at x = 0 the displacement $y$(x, t) is observed to decrease from its maximum value to half of its maximum value in time t = 2 $\times$ 10$^{-3}$ s.
\begin{qoptions}
\opt{A}{$y$(x, t) represents the motion of travelling pulse moving along positive X direction}
\opt{B}{$y$(x, t) does not represent the motion of a travelling pulse}
\opt{C}{Speed of wave pulse is 20 m/s}
\opt{D}{Speed of wave pulse is 10 m/s}
\end{qoptions}
\qn{14} There are three sound sources of frequencies 600Hz, 603Hz and 609Hz sounded simultaneously. Initially they are in same phase. A detector is placed close to them. Let I$_0$ be the intensity of each sound source. If detector is able to detect intensity nearly equal to
\begin{qoptions}
\opt{A}{3I$_0$, the frequency detected by the detector is 9Hz}
\opt{B}{3I$_0$, the frequency detected by the detector is 6Hz}
\opt{C}{9I$_0$, the frequency detected by the detector is 3Hz}
\opt{D}{9I$_0$, the frequency detected by the detector is 6Hz}
\end{qoptions}
\paraq{15}{16}A string 120 cm in length sustains a standing wave with the points of the string at which the displacement amplitude is equal to 3.5 mm being separated by 15.0 cm. Find
\qn{15} The maximum displacement amplitude
\begin{qoptions}
\opt{A}{$\sqrt2$ mm}
\opt{B}{2.95 mm}
\opt{C}{4.95 mm}
\opt{D}{6.95 mm}
\end{qoptions}
\qn{16} To which overtone do these oscillations correspond ?
\begin{qoptions}
\opt{A}{4}
\opt{B}{3}
\opt{C}{5}
\opt{D}{6}
\end{qoptions}
\paraq{17}{18}In the figure shown a sinusodial wave is generated at the end `A'. The wave travels along positive x-axis and during its motion it encounter another string BC at the junction `B' at x = 0. The density of strings AB and BC are $\rho$ and 9$\rho$ respectively and their radii of cross sections are 2r and r. The wave function, amplitude and wavelength of incident wave are respectively y$_\mathrm i$, A$_\mathrm i$ and $\lambda_\mathrm t$. Similarly for reflected and transmitted wave these parameters are y$_\mathrm r$, A$_\mathrm r$, $\lambda_\mathrm r$ and y$_\mathrm t$, A$_\mathrm t$, $\lambda_\mathrm t$ respectively.\par
\qfig{p1_phy_s0_q16_1}
\qn{17} The ratio of wavelengths $\lambda_\mathrm r$ to $\lambda_\mathrm t$ (i.e. $\lambda_\mathrm r:\lambda_\mathrm t$) will be
\begin{qoptions}
\opt{A}{1 : 1}
\opt{B}{3 : 2}
\opt{C}{2 : 3}
\opt{D}{None of these}
\end{qoptions}
\qn{18} The ratio of amplitudes A$_\mathrm r$ to A$_\mathrm t$ is (i.e. A$_\mathrm r$ : A$_\mathrm t$) will be
\begin{qoptions}
\opt{A}{1 : 1}
\opt{B}{1 : 4}
\opt{C}{4 : 1}
\opt{D}{None of these}
\end{qoptions}
\raggedcolumns
\papersec{1}
\partsec{2}{CHEMISTRY}
\qn{19} The solubility of ethylamine in a solution with pH = 8.5 is (K$_\mathrm{sp}$ of ethylamine = 1.78 $\times$ 10$^{-10}$)\par
\textbf{(Given:} $\sqrt{10}=3.16$ \textbf{)}
\begin{qoptions}
\opt{A}{5.6 $\times$ 10$^{-3}$ M}
\opt{B}{3.8 $\times$ 10$^{-4}$ M}
\opt{C}{5.6 $\times$ 10$^{-5}$ M}
\opt{D}{3.16 $\times$ 10$^{-6}$ M}
\end{qoptions}
\qn{20} A 250 ml flask containing NO$_\mathrm{(g)}$ at 0.42 atm is connected to a 100 ml flask containing oxygen gas at 0.84 atm by means of a stop cock at 350 K. The gases are mixed by opening the stop cock where the following equilibrium is established.\par
2NO(g) + O$_2$(g)$\longrightarrow$2NO$_2$(g)$\rightleftharpoons$N$_2$O$_4$(g)\par
The first reaction is complete while the second is at equilibrium. If the total pressure is 0.365 atm, K$_\mathrm p$ for dimersation of NO$_2$ is
\begin{qoptions}
\opt{A}{10 atm$^{-1}$}
\opt{B}{0.4 atm$^{-2}$}
\opt{C}{0.16 atm$^{-2}$}
\opt{D}{0.4 atm}
\end{qoptions}
\qn{21} In the excess of NH$_3$ (aq.), Cu$^{2+}$ ions forms a deep blue complex ion [Cu(NH$_3$)$_4$]$^{2+}$ which has K$_\mathrm f$ = 10$^{11}$ the concentration of Cu$^{2+}$ ion in solution prepared by adding 5 $\times$ 10$^{-3}$ mole of CuSO$_4$ to 0.5 L of 0.4 M NH$_3$ (aq) is
\begin{qoptions}
\opt{A}{$\dfrac{5\times10^{-5}}{(36)^4}$}
\opt{B}{$\dfrac{10^{-6}}{(36)^4}$}
\opt{C}{$\dfrac{10^{-5}}{(36)^4}$}
\opt{D}{$\dfrac{10^{-2}}{(39)^4}$}
\end{qoptions}
\qn{22} Equilibrium constants for the following reactions are given (\textbf{in atm}) for the reactions at 0$^\circ$C. Select the option mentioning correct order of \textbf{True (T)} or \textbf{False (F)} statement.\par
A.6H$_2$O$_{(s)}$ $\rightleftharpoons$ A.2H$_2$O$_{(s)}$ + 4H$_2$O$_{(g)}$ \hfill K$_\mathrm p$ = 1.6 $\times$ 10$^{-11}$\par
B.12H$_2$O$_{(s)}$ $\rightleftharpoons$ B.7H$_2$O$_{(s)}$ + 5H$_2$O$_{(g)}$ \hfill K$_\mathrm p$ = 2.43 $\times$ 10$^{-13}$\par
C.10H$_2$O$_{(s)}$ $\rightleftharpoons$ C$_{(s)}$ + 10H$_2$O$_{(g)}$ \hfill K$_\mathrm p$ = 10$^{-30}$\par
Aqueous tension of H$_2$O at 0$^\circ$C is given as 0.76 torr\par
(I) The most effective drying agent will be C$_{(s)}$ , out of C$_{(s)}$ , B.7H$_2$O$_{(s)}$ \& A.2H$_2$O$_{(s)}$\par
(II) At 0$^\circ$C A.6H$_2$O$_{(s)}$ \& B.12H$_2$O$_{(s)}$ will always be efflorescent.\par
(III) If R.H is less than 100\% in a chamber at 0$^\circ$C then none of the substance can act as deliquescent.
\begin{qoptions}
\opt{A}{TTT}
\opt{B}{FTF}
\opt{C}{TTF}
\opt{D}{TFF}
\end{qoptions}
\qn{23} The indicator constant for an acidic indicator, HIn, is 5$\times$10$^{-6}$ M. This indicator appears only in the colour of acidic form when $\dfrac{[\mathrm{In}^-]}{[\mathrm{HIn}]}\le\dfrac1{20}$ and it appears only in the colour of basic form when $\dfrac{[\mathrm{HIn}]}{[\mathrm{In}^-]}\le\dfrac1{40}$. The pH range of indicator is
\begin{qoptions}
\opt{A}{4.3 -- 6.3}
\opt{B}{4.0 -- 6.6}
\opt{C}{4.0 -- 6.9}
\opt{D}{3.7 -- 6.6}
\end{qoptions}
\qn{24} For the equilibrium CH$_3$--CH$_2$--CH$_2$--CH$_3$(g) $\rightleftharpoons$ CH$_3$--$\underset{\mathrm H}{\overset{\mathrm{CH}_3}{\mathrm C}}$--CH$_3$(g) equilibrium constant is found to be 1.732 at 298 K. Now if in a vessel at 298 K, a mixture of these two gases be taken as represented by the point P in the figure, predict what will happen :\par
\qfig{p1_chem_s0_q24_1}
\begin{qoptions}
\opt{A}{Immediately, above equilibrium will be set up}
\opt{B}{Above reaction will be in the forward direction till it attains equilibrium}
\opt{C}{Above reaction will go in the backward direction till it attains equilibrium}
\opt{D}{Nothing can be said}
\end{qoptions}
\qn{25} \qfig{p1_chem_s0_q25_1}
If valves are opened sequentially (A to D) \& each valve lets out 20 mL of stock solution, then select correct statements (Temp = 25$^\circ$C)
\begin{qoptions}
\opt{A}{After opening (A), buffer is created in beaker "P".}
\opt{B}{After opening (B), buffer is created beaker "Q".}
\opt{C}{After opening (C), contents in beaker Q can still act as buffer even though solutions pH decreases.}
\opt{D}{After opening (D), solution of pH $<$ 7 is created in beaker "R".}
\end{qoptions}
\qn{26} If two gases AB$_2$ and B$_2$C are mixed, following equilibria are readily established :\par
AB$_2$(g) + B$_2$C(g)$\rightleftharpoons$AB$_3$(g) + BC(g);\par
BC(g) + B$_2$C(g)$\rightleftharpoons$B$_3$C$_2$(g)\par
If the reaction is started only with AB$_2$ with B$_2$C, then which of the following is necessarily true at equilibrium ?
\begin{qoptions}
\opt{A}{[AB$_3$]$_\mathrm{eq}$ = [BC]$_\mathrm{eq}$}
\opt{B}{[AB$_2$]$_\mathrm{eq}$ $>$ [B$_2$C]$_\mathrm{eq}$}
\opt{C}{[AB$_3$]$_\mathrm{eq}$ $>$ [B$_3$C$_2$]$_\mathrm{eq}$}
\opt{D}{[AB$_3$]$_\mathrm{eq}$ $>$ [BC]$_\mathrm{eq}$}
\end{qoptions}
\qn{27} A solution contains 0.05 M each of NaCl and Na$_2$CrO$_4$. Solid AgNO$_3$ is gradually added to it. Which of the following facts is/are incorrect?\par
Given: K$_\mathrm{sp}$(AgCl) = 1.7 x 10$^{-10}$ M$^2$ and K$_\mathrm{sp}$(Ag$_2$CrO$_4$) = 1.9 x 10$^{-12}$ M$^3$.
\begin{qoptions}
\opt{A}{Cl$^-$ ions are precipitated first}
\opt{B}{CrO$_4^{2-}$ ions are precipitated first}
\opt{C}{Both Cl$^-$ and CrO$_4^{2-}$ ions are precipitated simultaneously}
\opt{D}{The second ion starts precipitating when half of the first ion is precipitated}
\end{qoptions}
\qn{28} CN$^-_{\mathrm{(aq.)}}$ + H$_2$O $\rightleftharpoons$ HCN$_{\mathrm{(aq.)}}$ + OH$^-_{\mathrm{(aq.)}}$\par
Addition of which of the following reagent can make the reaction proceed towards forward direction.
\begin{qoptions}
\opt{A}{H$_2$O}
\opt{B}{HCl}
\opt{C}{NaBr}
\opt{D}{KOH}
\end{qoptions}
\qn{29} 1 M solution of a dibasic acid H$_2$A $\left(\mathrm K_{\mathrm a_1}=10^{-6}\text{ and }\mathrm K_{\mathrm a_2}=10^{-10}\right)$ is titrated against 1 M NaOH solution. Choose the \textbf{correct} statements regarding the relative concentration of H$_2$A, HA$^-$ and A$^{2-}$ at various pH during titration.
\begin{qoptions}
\opt{A}{At pH = 3 ; [H$_2$A] is maximum.}
\opt{B}{At pH = 6.5 ; [HA$^-$] is maximum.}
\opt{C}{At pH = 10 ; [H$_2$A] is minimum.}
\opt{D}{At pH = 8 ; [A$^{2-}$] is maximum.}
\end{qoptions}
\qn{30} For the reaction N$_2$O$_5$(g) $\longrightarrow$ 2NO$_2$(g) + $\dfrac12$O$_2$(g)\par
N$_2$O$_5$(g) decomposed at a constant volume and temperature. Initial pressure of N$_2$O$_5$ is 600 mm of Hg and total pressure at equilibrium is 960 mm of Hg. Assuming ideal behavior of gases, Select the correct option(s).
\begin{qoptions}
\opt{A}{Involved reaction is intramolecular redox reaction}
\opt{B}{The mole fraction of NO$_2$ in the equilibrium mixture is 0.5}
\opt{C}{The equilibrium constant K$_\mathrm p$ for the above reaction is 540 mm of Hg}
\opt{D}{For the above reaction K$_\mathrm C$ = K$_\mathrm p$ (RT)$^3$}
\end{qoptions}
\qn{31} Solubility product of AgCl (s) in water at 25$^\circ$C is \ 12 $\times$ 10$^{-10}$ M$^2$ then select the correct options.
\begin{qoptions}
\opt{A}{Solubility of AgCl in presence of \ 10$^{-5}$ M AgNO$_3$ is \ 3 $\times$ 10$^{-5}$ moles/L.}
\opt{B}{Solubility of AgCl in presence of excess of KCl will increase.}
\opt{C}{Solubility of AgCl in presence of 2M NH$_3$ is greater than that in pure water.}
\opt{D}{Solubility of AgCl remains unaffected on adding any amount of NaBr because there is no common ion effect.}
\end{qoptions}
\qn{32} Ammonia under a pressure of 20 atm at 127$^\circ$C is heated to 327$^\circ$C in a closed rigid vessel. Under these conditions, ammonia is partially decomposed to N$_2$ and H$_2$ according to the reaction\par
2NH$_3$(g) $\rightleftharpoons$ N$_2$(g) \ + 3H$_2$(g)\par
After decomposition, the pressure inceases to 50 atm. Which of the following statement(s) is/are \textbf{correct} at 327$^\circ$C?
\begin{qoptions}
\opt{A}{The degree of dissociation of NH$_3$ is 2/3}
\opt{B}{The K$_\mathrm p$ of the reaction is 2.7 $\times$ 10$^3$atm$^2$.}
\opt{C}{The partial pressures of both, N$_2$ and NH$_3$ gas at equilibrium is 10 atm.}
\opt{D}{The pressure of NH$_3$ at equilibrium is less than that of N$_2$.}
\end{qoptions}
\paraq{33}{34}Consider the following experimental set-up with stop cocks: S$_1$, S$_2$, S$_3$, S$_4$ and taps : T$_1$, T$_2$, T$_3$, T$_4$\par
\qfig{p1_chem_s0_q32_1}
\qn{33} In which of the following cases will the solution in tub have a pH different from others:
\begin{qoptions}
\opt{A}{T$_1$ and T$_4$ are opened}
\opt{B}{only T$_2$ is opened}
\opt{C}{T$_1$, T$_2$, T$_3$ and T$_4$ are all opened}
\opt{D}{T$_1$ is opened and the content of bucket is poured in the tub.}
\end{qoptions}
\qn{34} A pH meter is immersed in the bucket of the original set up shown. It shows a reading of 7. Now S$_1$ is the first opened and all content is allowed to drain out into the bucket. pH of the solution in the bucket is measured again. Similarly S$_2$, S$_3$ and S$_4$ are sequentially opened and the pH of the solution in the bucket is measured each time after complete drainage. Opening which stop cock would cause maximum pH change (increase/decrease) compared to previous pH in bucket.
\begin{qoptions}
\opt{A}{S$_1$}
\opt{B}{S$_2$}
\opt{C}{S$_3$}
\opt{D}{S$_4$}
\end{qoptions}
\paraq{35}{36}Following reaction is at equilibrium in basic medium at 300 K.\par
I$_2$ (s) \ $\rightleftharpoons$ I$^-$ (aq.) + IO$_3^-$ (aq.)\par
Equilibrium concentration at 300 K are , [ I$^-$] = 0.1 M, [IO$_3^-$ ] = 0.1 M\par
Given $\Delta$G$^\circ_\mathrm f$ (I$^-$, aq) = $-$ 50 kJ/ mol\par
$\Delta$G$^\circ_\mathrm f$ (IO$_3^-$, aq) = $-$ 123.5 kJ/ mol\par
$\Delta$G$^\circ_\mathrm f$ (H$_2$O, $l$ ) = $-$ 233 kJ/ mol\par
$\Delta$G$^\circ_\mathrm f$ (OH$^-$, aq ) = $-$ 150 kJ/ mol\par
R = $\dfrac{25}3$ J mol$^{-1}$ K$^{-1}$ \ log$_{10}$ e = 2.3
\qn{35} Value of equilibrium constant is
\begin{qoptions}
\opt{A}{10$^{25}$}
\opt{B}{10$^{35}$}
\opt{C}{10$^{30}$}
\opt{D}{10$^{27}$}
\end{qoptions}
\qn{36} pH value at equilibrium is
\begin{qoptions}
\opt{A}{8}
\opt{B}{6}
\opt{C}{9}
\opt{D}{5}
\end{qoptions}
\raggedcolumns
\papersec{1}
\partsec{3}{MATHEMATICS}
\qn{37} Let P(z) = z$^3$ + az$^2$ + bz + c, where a, b, c, z $\in$ R. There exist a complex number $\omega$ such that the three roots of P(z) are 2$\omega$ $-$ 4, $\omega$ + 9i, $\omega$ + 3i, where i$^2$ = $-$1. Then value of $|\mathrm a+\mathrm b+\mathrm c|$ is equal to :
\begin{qoptions}
\opt{A}{121}
\opt{B}{136}
\opt{C}{79}
\opt{D}{does not exist}
\end{qoptions}
\qn{38} The polynomial $\mathrm p(x)=(1+x+x^2+\dots\dots+x^{17})^2-x^{17}$ has 34 complex roots of the form $z_k=r_k[\cos(2\pi a_k)+i\sin(2\pi a_k)]$, k = 1, 2, 3.... 34 with $0<a_1\le a_2\le a_3\le\dots..\le a_{34}<1$ and $r_k>0$, then value of $a_1+a_2+a_3+a_4+a_5$ is equal to :
\begin{qoptions}
\opt{A}{$\dfrac1{19}$}
\opt{B}{$\dfrac3{19}$}
\opt{C}{$\dfrac{159}{323}$}
\opt{D}{$\dfrac{482}{591}$}
\end{qoptions}
\qn{39} Complex numbers a, b and c are zeros of p(z) = z$^3$ + qz + r, where q, r $\in$ R and $|\mathrm a|^2+|\mathrm b|^2+|\mathrm c|^2=250$. If the points corresponding to a, b and c are in the complex plane are the vertices of right angled triangle with hypotenuse `h' then $\left(\dfrac{\mathrm h^2}{125}\right)$ is equal to :
\begin{qoptions}
\opt{A}{1}
\opt{B}{2}
\opt{C}{3}
\opt{D}{4}
\end{qoptions}
\qn{40} Coefficient of $x^2$ in\par$\left(\underbrace{\dots\dots(x(x-2)^2-2)^2-\dots\dots-2}_{100\text{ times}}\right)^2$\par is equal to :
\begin{qoptions}
\opt{A}{$\dfrac{4^{199}-4^{99}}3$}
\opt{B}{$\dfrac{4^{201}-4^{100}}3$}
\opt{C}{$\dfrac{4^{99}-4^{50}}3$}
\opt{D}{None of these}
\end{qoptions}
\qn{41} If $\displaystyle\sum_{r=1}^{2008}(2009-r)\cos\left(\dfrac{2\pi r}{2009}\right)=-\dfrac n2$, then n is equal to :
\begin{qoptions}
\opt{A}{2008}
\opt{B}{2009}
\opt{C}{2016}
\opt{D}{2018}
\end{qoptions}
\qn{42} The coefficients of x$^n$ in the polynomial\par$\left(x+{}^{2n+1}\mathrm C_0\right)\left(x+{}^{2n+1}\mathrm C_1\right)\left(x+{}^{2n+1}\mathrm C_2\right)\dots\dots$\par$\left(x+{}^{2n+1}\mathrm C_n\right)$ :
\begin{qoptions}
\opt{A}{$2^{n+1}$}
\opt{B}{$2^{2n+1}-1$}
\opt{C}{$2^{2n}$}
\opt{D}{None of these}
\end{qoptions}
\qn{43} Let $z_1$ and $z_2$ are two complex numbers such that $z_1+z_2=z_1^2+z_2^2=\dfrac{2\mathrm i}{\sqrt3}$ (where $\mathrm i=\sqrt{-1}$) then :
\begin{qoptions}
\opt{A}{$z_1$ and $z_2$ lies in first quadrant}
\opt{B}{$z_1$ and $z_2$ lies in second quadrant}
\opt{C}{$z_1$ and $z_2$ lies in third quadrant}
\opt{D}{$z_1$ and $z_2$ lies in different quadrant}
\end{qoptions}
\qn{44} If S$_\mathrm n$ = 2($\omega$+1)($\omega^2$+1) +3(2$\omega$+1)(2$\omega^2$+1) + 4(3$\omega$+1)(3$\omega^2$+1)+..... upto n terms then the correct statement(s) is are : (where n $\ge$ 3 and $\omega$ is non real cube root of unity)
\begin{qoptions}
\opt{A}{S$_\mathrm n$ can not be a prime number}
\opt{B}{unit digit of S$_{24}$ is 4}
\opt{C}{S$_\mathrm n$ is always divisible by n}
\opt{D}{S$_\mathrm n$ is an increasingh sequence}
\end{qoptions}
\qn{45} For two non-zero complex numbers A and B if $\mathrm A+\dfrac1{\mathrm B}=\bar{\mathrm A}$ and $\dfrac1{\mathrm A}+\mathrm B=\bar{\mathrm B}$, then :
\begin{qoptions}
\opt{A}{A and B are purely imaginary}
\opt{B}{A and B are purely real}
\opt{C}{AB=$-\dfrac12$}
\opt{D}{AB=$\dfrac12$}
\end{qoptions}
\qn{46} If 1, $\omega$, $\omega^2$ .....$\omega^{n-1}$ are the n$^\mathrm{th}$ roots of unity, then the value of $(4-\omega)(4-\omega^2)\dots\dots\dots(4-\omega^{n-1})$ is equal to :
\begin{qoptions}
\opt{A}{$\dfrac{2^{2n}-1}3$}
\opt{B}{can not be prime number}
\opt{C}{${}^n\mathrm C_1+{}^n\mathrm C_2\,3+\dots\dots+{}^n\mathrm C_n\,3^{n-1}$}
\opt{D}{$\dfrac{6^n+1}5$}
\end{qoptions}
\qn{47} Let $z_1$ and $z_2$ are two distinct non-real complex numbers such that $z_1+\dfrac1{z_2}\in\mathrm R$ and $z_2+\dfrac1{z_1}\in\mathrm R$, then value of $|z_1z_2|$ may
\begin{qoptions}
\opt{A}{A prime number}
\opt{B}{A natural number}
\opt{C}{Positive integer}
\opt{D}{Negative integer}
\end{qoptions}
\qn{48} If $(1+x)^n=c_0+c_1x+c_2x^2+\dots\dots.+c_nx^n$, where n is a positive integer, then
\begin{qoptions}
\opt{A}{$c_0-c_2+c_4-\dots\dots\dots=2^{n/2}\cos\left(\dfrac{n\pi}4\right)$}
\opt{B}{$c_1-c_3+c_5-\dots\dots\dots=2^{n/2}\sin\dfrac{n\pi}4$}
\opt{C}{$c_0+c_4+c_8+\dots\dots\dots=2^{n-2}+2^{\frac{n-2}2}\cos\left(\dfrac{n\pi}4\right)$}
\opt{D}{$c_0+c_3+c_6+\dots\dots\dots=\dfrac13\left(2^n+2\cos\dfrac{n\pi}3\right)$}
\end{qoptions}
\qn{49} Let $a_n=\dfrac{(1000)^n}{n!}$ for n $\in$ N. Then a$_n$ is greatest, when
\begin{qoptions}
\opt{A}{n = 998}
\opt{B}{n = 999}
\opt{C}{n = 1000}
\opt{D}{n = 1001}
\end{qoptions}
\qn{50} If $(1+2x+3x^2)^{10}=\mathrm a_0+\mathrm a_1+\mathrm a_2x^2+\dots\dots..+\mathrm a_{20}x^{20}$ then
\begin{qoptions}
\opt{A}{a$_1$ = 20}
\opt{B}{a$_2$ = 210}
\opt{C}{a$_4$ = 8085}
\opt{D}{a$_{20}$ = 2$^2\cdot3^7\cdot7$}
\end{qoptions}
\paraq{51}{52}Let c$_1$ and c$_2$ be the two curves on the complex plane defined as\par
c$_1$ : $z+\bar z=2|z-1|$ \ and \ c$_2$ : arg$(z+1+i)=\alpha$, where $\alpha\in(0,\pi)$,\par
such that curves c$_1$ and c$_2$ \ touch each other at p(z$_1$).
\qn{51} A particle starts from a point P(z$_0$). It moves horizontally away from origin by 2 units and then vertically away from origin by 3 units to reach at a point Q(z$_1$). Then $|z_1|$ is equal to :
\begin{qoptions}
\opt{A}{3}
\opt{B}{5}
\opt{C}{7}
\opt{D}{9}
\end{qoptions}
\qn{52} If P(z$_0$) is rotated about origin through an angle 2$\alpha$ in clockwise sense, then the area bounded by c$_1$, and the line joining P(z$_0$) and Q(z$'_0$) is, (where Q(z$'_0$)) is the new position of P(z$_0$) after rotation :
\begin{qoptions}
\opt{A}{$\dfrac56$ sq. unit}
\opt{B}{1 sq. unit}
\opt{C}{2 sq. unit}
\opt{D}{$\dfrac23$ sq. unit}
\end{qoptions}
\paraq{53}{54}$\mathrm e^x=1+x+\dfrac{x^2}{2!}+\dfrac{x^3}{3!}+\dots\dots.\infty$\par
$f(x)=f(0)+f'(0)x+f''(0)\dfrac{x^2}{2!}+\dots\dots\dots\infty$
\qn{53} If $\dfrac{\mathrm e^x}{1-x}=\mathrm B_0+\mathrm B_1x+\mathrm B_2x^2+\dots\dots..+\mathrm B_nx^n+\dots\dots$ then the value of B$_\mathrm n$ $-$ B$_{\mathrm n-1}$ is equal to :
\begin{qoptions}
\opt{A}{n!}
\opt{B}{(n + 1)!}
\opt{C}{$\dfrac1{\mathrm n!}$}
\opt{D}{$\dfrac1{(\mathrm n+1)!}$}
\end{qoptions}
\qn{54} Sum of the series $\left(\dfrac2{3!}+\dfrac4{5!}+\dfrac6{7!}+\dots\dots.\infty\right)$
\begin{qoptions}
\opt{A}{e$^2$}
\opt{B}{e$^{-1}$}
\opt{C}{e$^{-2}$}
\opt{D}{None of these}
\end{qoptions}
\end{multicols}
\end{document}
